% Contenido para: sections/4experiencia.tex

La ejecución de la experiencia de laboratorio se llevó a cabo siguiendo rigurosamente los lineamientos establecidos en la guía, buscando una comprensión práctica del osciloscopio digital. Para ello, se utilizaron dos instrumentos principales: un Generador de Señales Rigol DG1022 (2 Canales, 20 MHz, 100 MSa/s; N/S: 010000029711) para generar las señales de prueba, y un Osciloscopio Digital Rigol DS4014 (4 Canales, 100 MHz, 4 GS/s; N/S: 010000012635) para su visualización y análisis. Durante todo el proceso, se mantuvo la precaución de evitar el uso de la función de auto-configuración ("Auto-setting") del osciloscopio, a fin de asegurar una manipulación manual y consciente de todos los parámetros.

\subsection{Análisis del Sistema de Disparo (Trigger y Nivel)}
El estudio del sistema de disparo, o trigger, comienza con la conexión directa de la salida del generador de funciones al Canal 1 (CH1) del osciloscopio. Se aplica una señal sinusoidal de \SI{50}{\hertz} y se procede a ajustar manualmente la ganancia vertical (V/div) y la base de tiempo (s/div) para obtener una visualización clara de la traza. El propósito de esta fase es comprender cómo el trigger es esencial para sincronizar el barrido horizontal con la señal de entrada, logrando así una imagen estable en la pantalla.

Inicialmente, el trigger se configura para sincronizar con CH1, estableciendo un nivel de disparo (Level) ligeramente por encima de cero (\SI{0}{\volt}$^+$) y seleccionando el flanco de subida. Para entender el impacto del nivel de disparo en la estabilidad de la traza, este se varía progresivamente desde \SI{0}{\volt}$^+$ hacia valores positivos mayores, y luego se repite el proceso variando el nivel desde \SI{0}{\volt}$^+$ hacia valores negativos. El objetivo es determinar experimentalmente el rango de niveles válidos para lograr una sincronización efectiva con la señal aplicada.

Posteriormente, se investiga el efecto del flanco de disparo. Manteniendo la sincronización con CH1 y el nivel cercano a cero, se conmuta entre el flanco de pendiente positiva (subida) y negativa (bajada). El propósito es observar cómo la elección del flanco afecta el punto de inicio de la visualización de la señal en la pantalla.

Para un análisis más complejo de la sincronización, se mantiene la señal sinusoidal en CH1 y se conecta una señal triangular de \SI{50}{\hertz} al Canal 2 (CH2). Se explora entonces la función de la fuente de trigger, alternando la selección entre CH1, CH2 y la sincronización con la red eléctrica (Line), con el fin de observar cómo afecta a la estabilidad de ambas señales la elección de la fuente de sincronización. Finalmente, para profundizar en esta relación, se modifica la frecuencia de la señal triangular en CH2 a múltiplos, submúltiplos y frecuencias no relacionadas con los \SI{50}{\hertz} de CH1. Este procedimiento se realiza para estudiar el comportamiento de la visualización cuando la fuente de trigger y una de las señales visualizadas tienen frecuencias distintas o no sincronizadas.

\subsection{Estudio de Acoplamientos y Modos de Visualización}
Esta sección se dedica a entender los distintos modos de acoplamiento de entrada del osciloscopio y las opciones de visualización de los canales. Se aplican dos señales sinusoidales, ambas con una componente de corriente continua (offset), una al CH1 y otra al CH2, utilizando el generador de funciones Rigol DG1022 para configurar los offsets de manera controlada.

Inicialmente, ambos canales se configuran en acoplamiento de \textbf{Corriente Continua (CC o DC)}. Este modo permite el paso de todas las componentes (AC y DC) de la señal. Se procede a medir la amplitud pico-pico, el valor medio y el período de cada señal, así como a calcular el desfase entre ellas. El objetivo es obtener una caracterización completa de las señales originales.

Posteriormente, se cambia el acoplamiento de ambos canales a \textbf{Corriente Alterna (CA o AC)}. Este modo filtra la componente de corriente continua. Se repite la medición del desfase y se registran las señales visualizadas, con el propósito de observar el efecto del filtrado de la componente DC. Para establecer una referencia clara de \SI{0}{\volt}, se selecciona brevemente el acoplamiento \textbf{GND (Tierra)} en ambos canales.

Para explorar las opciones de visualización, se vuelve al acoplamiento CC y se activa el modo \textbf{invertido} en ambos canales simultáneamente. Se realizan mediciones de amplitudes, valores medios y períodos para observar el efecto de la inversión. A continuación, se activa el modo invertido solo en CH1, manteniendo CH2 en modo normal. En esta configuración, se mide el desfase resultante entre CH1 y CH2, con el fin de determinar cómo la inversión de un solo canal altera la relación de fase medida.

\subsection{Análisis de Ganancias, Base de Tiempo y Modos Temporales}
Esta fase se centra en la optimización de la visualización mediante los ajustes de escala vertical (ganancia) y horizontal (base de tiempo), así como en la exploración de los modos Y-t y X-Y. El osciloscopio se configura en el modo \textbf{Y-t}. Se aplica una señal sinusoidal de \SI{5}{\volt_{rms}} y \SI{50}{\hertz} al CH1, y otra señal sinusoidal de \SI{1}{\volt_{rms}} y \SI{50}{\hertz} al CH2, ambas generadas por el Rigol DG1022.

El primer paso es ajustar las ganancias verticales (V/div) y la base de tiempo (s/div) para visualizar claramente al menos dos ciclos completos de ambas señales. Durante este ajuste, se realizan mediciones de amplitud y frecuencia. A continuación, se exploran los límites de la visualización variando la ganancia: se reduce progresivamente la ganancia de un canal hasta que la señal deja de ser apreciable, y luego se aumenta progresivamente hasta que la señal se satura visiblemente en la pantalla. El objetivo de este procedimiento es identificar el rango de ajustes de ganancia que permiten una medición válida y evitar condiciones de sub-visualización o saturación.

Posteriormente, el osciloscopio se cambia al modo \textbf{X-Y}. Con las señales de \SI{50}{\hertz} inicialmente configuradas en ambos canales, se registra la figura de Lissajous resultante y se miden sus extensiones máximas en los ejes X e Y. Para estudiar la relación entre la fase y la forma de estas figuras, se introduce un desfase variable en una de las señales, registrando los cambios en la figura. Finalmente, para analizar el efecto de la relación de frecuencias, se modifica la frecuencia de una de las señales (manteniendo la otra en \SI{50}{\hertz}) y se registran las figuras de Lissajous correspondientes a diferentes relaciones de frecuencia (ej. múltiplos enteros).

\subsection{Medición de Señal Transitoria}
La experiencia concluye con la configuración del osciloscopio para la captura de una señal transitoria no periódica, inyectada desde la mesa central del laboratorio. Para lograr su registro, se ajustan los parámetros de trigger (seleccionando probablemente el modo Single para capturar un único evento), la ganancia vertical y la base de tiempo horizontal. Estos ajustes se realizan en función de las características temporales y de amplitud esperadas de la señal transitoria, con el objetivo de obtener una instantánea clara del evento para su posterior análisis.